\chapter[1948 --- Muller]{1948: Paul Hermann M\"uller}
\label{ch:1948_paulhermannmuller}

\section*{Main Contribution}

\textit{Discovered that DDT is a highly effective insecticide, enabling control of malaria and typhus---though later recognized as an environmental hazard.}

\section{Official Citation}

for his discovery of the high efficiency of DDT as a contact poison against several arthropods

\section{Background and Historical Context}

Paul Hermann M\"uller was born on January 12, 1899, in Olten, Switzerland, into a middle-class family. His father was a postal employee, and his mother was a schoolteacher. M\"uller grew up in the small Swiss town of Solothurn, where he attended the local Gymnasium and demonstrated an early aptitude for chemistry and biology. After completing his secondary education, he enrolled at the University of Basel, where he studied chemistry and graduated with his Ph.D. in 1925 under the supervision of Paul Karrer, who would later win the Nobel Prize in Chemistry for his work on carotenoids and vitamins.

After completing his doctoral studies, M\"uller joined the research laboratories of J.\ R.\ Geigy AG, a chemical and pharmaceutical company based in Basel. This company, which would eventually become part of the multinational pharmaceutical giant Novartis, provided M\"uller with the resources and infrastructure to conduct the research that would lead to his Nobel Prize. His work at Geigy focused on the development of new chemical compounds with biological activity, particularly compounds that could be used as insecticides.

The scientific and public health context of the 1930s and 1940s was one of intense concern about insect-borne diseases. Malaria, typhus, and other vector-borne diseases were major causes of morbidity and mortality throughout the world, particularly in tropical and subtropical regions. The existing methods of insect control---including the use of pyrethrin, a natural insecticide derived from chrysanthemum flowers---were effective but expensive, difficult to produce in large quantities, and of limited duration. There was an urgent need for a cheap, effective, and long-lasting insecticide that could be used to control the spread of insect-borne diseases.

It was within this context of medical and scientific need that M\"uller made his discovery of DDT.

\section{The Discovery/Contribution}

Paul Hermann M\"uller received the Nobel Prize in Physiology or Medicine ``for his discovery of the high efficiency of DDT as a contact poison against several arthropods.'' His work demonstrated that dichlorodiphenyltrichloroethane (DDT), a synthetic organic compound, had remarkable insecticidal properties and could be used to control the insect vectors of major infectious diseases.

\subsection{The Discovery of DDT's Insecticidal Properties}

The discovery of DDT's insecticidal properties was, in a sense, the result of a systematic screening program rather than a single eureka moment. M\"uller and his colleagues at Geigy had been searching for new insecticidal compounds since the early 1930s, and they had developed an efficient method for testing candidate compounds for their ability to kill insects. The method involved spraying a solution of the test compound onto a surface and then placing insects on the treated surface to see whether they were killed.

In 1935, M\"uller tested a compound called dichlorodiphenyltrichloroethane---DDT---that had first been synthesized by Othmar Zeidler in 1874 but had never been tested for biological activity. The results were remarkable: DDT was lethal to a wide variety of insects, including houseflies, body lice, potato beetles, and mosquito larvae, at concentrations far lower than those required for other insecticides. Moreover, DDT was remarkably persistent, retaining its insecticidal activity for weeks or even months after application---a property that made it ideal for use in public health programs.

M\"uller published his findings in 1939, and the discovery was immediately recognized as a breakthrough of major significance. The Swiss government authorized the use of DDT for crop protection in 1940, and during World War II, the Allied forces used DDT extensively to control the insect vectors of malaria and typhus among troops operating in tropical and subtropical regions. The use of DDT is estimated to have saved millions of lives during the war and in the immediate postwar period.

\subsection{The Chemistry and Mechanism of Action of DDT}

DDT (1,1,1-trichloro-2,2-bis(4-chlorophenyl)ethane) is an organochlorine compound with the chemical formula C$_{14}$H$_9$Cl$_5$. It is a white, crystalline solid that is virtually insoluble in water but highly soluble in fats and oils. This lipophilicity (fat solubility) is responsible for DDT's persistence in the environment and its tendency to accumulate in the fatty tissues of organisms.

DDT acts as a contact insecticide, meaning that it kills insects upon direct contact with their body surface. The mechanism of action of DDT involves the disruption of sodium ion channels in the nerve cell membrane. Normally, sodium channels open briefly during the generation of an action potential, allowing sodium ions to flow into the cell and then quickly closing to allow the cell to repolarize. DDT binds to the sodium channel and prevents it from closing properly, causing a continuous influx of sodium ions that leads to repetitive firing of nerve impulses, convulsions, paralysis, and ultimately death.

This mechanism of action---the disruption of voltage-gated sodium channels---has since been found to be shared by several other classes of insecticides, including pyrethroids (synthetic analogs of the natural insecticide pyrethrin). The specificity of DDT for insect sodium channels, as opposed to mammalian sodium channels, is due to subtle differences in the structure of the channel proteins between insects and mammals.

\subsection{Field Trials and Large-Scale Application}

M\"uller's laboratory discovery was rapidly translated into practical applications through a series of field trials and large-scale deployment programs. In 1943, during a typhus epidemic in Naples, Italy, the Allied forces used DDT to delouse the civilian population. The compound was applied as a dust to the clothing and bodies of thousands of people, and the epidemic was brought under control within weeks. This success demonstrated the potential of DDT for the control of insect-borne diseases and led to its widespread adoption by public health agencies around the world.

After World War II, DDT was used extensively in tropical and subtropical regions to control malaria. The World Health Organization (WHO) launched a global malaria eradication program in 1955 that relied heavily on DDT for indoor residual spraying---the application of DDT to the interior walls of homes, where it kills mosquitoes that land on the treated surfaces. This program was remarkably successful in many regions, reducing malaria incidence by over 90\% in countries such as India, Sri Lanka, and Greece.

\section{Impact on Modern Medicine}

M\"uller's discovery of DDT has had a complex and multifaceted impact on modern medicine, encompassing both the enormous benefits of insect-borne disease control and the environmental and health consequences of DDT's widespread use.

\subsection{Insect-Borne Disease Control}

The most direct medical impact of M\"uller's discovery was the control of insect-borne diseases. Before the introduction of DDT, malaria was one of the leading causes of death in the world, responsible for an estimated 1 to 2 million deaths per year. The use of DDT for indoor residual spraying contributed to a dramatic reduction in malaria incidence in many parts of the world during the 1950s and 1960s.

Typhus, which is transmitted by body lice, was also effectively controlled by DDT. During World War II, the use of DDT to delouse military personnel and civilians prevented large-scale typhus epidemics that had been a major source of mortality in previous wars. The control of typhus was one of the major public health achievements of the war and demonstrated the potential of chemical insecticides to prevent epidemic disease.

Other insect-borne diseases that were affected by DDT control programs include Chagas disease (transmitted by triatomine bugs), leishmaniasis (transmitted by sandflies), and sleeping sickness (transmitted by tsetse flies). While DDT was not always effective against all of these vectors, its use contributed to significant reductions in the burden of these diseases in many regions.

\subsection{The Environmental Revolution and the Banning of DDT}

Despite its enormous benefits, DDT's widespread use also had serious environmental consequences. Because DDT is highly persistent and lipophilic, it accumulates in the fatty tissues of organisms and biomagnifies through the food chain. Top predators---including birds of prey, fish, and marine mammals---accumulate the highest concentrations of DDT, leading to a range of toxic effects including eggshell thinning in birds, reproductive failure, and immune system dysfunction.

Rachel Carson's landmark 1962 book \textit{Silent Spring} brought these environmental consequences to the attention of the general public and helped spark the modern environmental movement. Carson documented the devastating effects of DDT on bird populations and other wildlife, and she argued for the development of more targeted and environmentally friendly methods of insect control. The book was initially met with fierce opposition from the chemical industry, but it ultimately led to a fundamental reassessment of the environmental impact of pesticides.

In 1972, the United States Environmental Protection Agency (EPA) banned the use of DDT in the United States, and in 2001, the Stockholm Convention on Persistent Organic Pollutants restricted the use of DDT to disease vector control in countries where malaria remains a major public health threat. This compromise reflects the tension between the enormous benefits of DDT for disease control and its significant environmental costs.

\subsection{Development of Safer Insecticides}

The environmental concerns raised by DDT spurred the development of alternative insecticides with more favorable environmental profiles. Pyrethroids, which share DDT's mechanism of action but are less persistent and less bioaccumulative, have become the insecticides of choice for many public health applications. Other classes of insecticides, including organophosphates and carbamates, target different biochemical pathways and offer alternative strategies for insect control.

The development of integrated pest management (IPM) strategies---which combine chemical, biological, and environmental methods of insect control---has also been driven, in part, by the lessons learned from DDT's environmental impact. IPM approaches seek to minimize the use of chemical insecticides while maintaining effective control of insect-borne diseases, representing a more sustainable approach to vector control.

\subsection{Ongoing Use in Malaria Control}

Despite the environmental concerns, DDT continues to be used for indoor residual spraying in countries where malaria remains endemic and where alternative insecticides are too expensive or unavailable. The WHO currently recommends DDT as one of several options for indoor residual spraying, and approximately 10 to 15 countries, primarily in sub-Saharan Africa and South Asia, continue to use DDT for malaria control. The ongoing use of DDT in these countries reflects the continuing importance of M\"uller's discovery for global public health, even as efforts are made to transition to safer alternatives.

\section{Legacy}

Paul Hermann M\"uller's legacy is complex and multifaceted, reflecting both the enormous benefits and the significant costs of his discovery. The identification of DDT as an effective insecticide was one of the major public health achievements of the twentieth century, and it contributed directly to the control of malaria, typhus, and other insect-borne diseases that had claimed millions of lives. M\"uller's work demonstrated the power of systematic chemical research to address pressing medical problems and led to the development of an entire field of insecticide chemistry.

At the same time, the environmental consequences of DDT's widespread use have served as a cautionary tale about the importance of considering the ecological impact of technological innovations. The story of DDT has influenced the development of environmental regulations and has helped shape public attitudes toward the use of chemicals in the environment.

M\"uller continued to work at the Geigy laboratories after receiving the Nobel Prize, contributing to the development of additional insecticidal compounds, including the organochlorine insecticide lindane and several organophosphorus compounds. He received numerous honors during his career, including honorary degrees from universities in several countries. He died on October 12, 1965, at the age of 66.

The story of Paul M\"uller and DDT is a reminder that scientific discoveries can have both positive and negative consequences, and that the task of balancing these consequences is one of the great challenges of modern society. As we face new environmental and public health challenges---from climate change to emerging infectious diseases---the lessons of DDT continue to inform our thinking about the relationship between technology, health, and the environment.


\section{Modern Developments}

Müller's DDT work has led to a complex legacy. DDT was banned in many countries due to environmental persistence and bioaccumulation, but remains used for malaria control in some regions. Understanding of endocrine disruption has led to regulations on pesticide use. New insecticides (e.g., chlorfenapyr, clothianidin) are being developed with improved safety profiles. Gene drive technology aims to suppress mosquito populations without chemical insecticides.


\section{Bibliography}

\begin{thebibliography}{9}
\bibitem{nobel-1948}
Nobel Prize Outreach, ``The Nobel Prize in Physiology or Medicine 1948,''\emph{NobelPrize.org}.\\
\url{https://www.nobelprize.org/prizes/medicine/1948/summary/}

\bibitem{lecture-1948}
Paul Hermann M"uller, ``Nobel Lecture,''\emph{NobelPrize.org}. Winner-authored primary source; downloadable from the official Nobel Prize archive.\\
\url{https://www.nobelprize.org/prizes/medicine/1948/muller/lecture/}
\end{thebibliography}
